\documentclass{beamer}
\usetheme{Madrid}
\usecolortheme{whale}

\title[FPGA NN Exploration]{Quantized and Gated Neural Network Accelerators on FPGA}
\subtitle{TTQ, Density-Adaptive Pruning, and Spatial Hysteresis}
\author{Research & Development Presentation}
\date{\today}

\begin{document}

% Slide 1: Title Slide
\begin{frame}
\titlepage
\end{frame}

% Slide 2: Work Done So Far
\begin{frame}{Work Done So Far: Accelerator Evolution}
We have implemented and evaluated a complete hardware-software co-design pipeline:
\begin{itemize}
    \item \textbf{Floating-Point Baseline}: High accuracy ($98.35\%$ MNIST, $90.62\%$ CIFAR-10) but heavy memory and arithmetic footprints.
    \item \textbf{Trained Ternary Quantization (TTQ)}:
    \begin{itemize}
        \item Weights compressed to 2-bits (16$\times$ storage reduction).
        \item Asymmetric learnable scaling factors ($W_p$, $W_n$).
        \item Replaced multiplications with ternary add/subtract (zero DSPs per tap).
    \end{itemize}
    \item \textbf{Activation Pruning Accelerators}:
    \begin{itemize}
        \item Fused Convolution and Pooling sequential pipeline.
        \item Lookahead skip-gating FSMs (2-tap for Conv, 4-tap for FC).
        \item Two gating algorithms: Density-Adaptive Activation Pruning (DAAP) and Spatial Hysteresis Gating.
    \end{itemize}
\end{itemize}
\end{frame}

% Slide 3: Detailed Thresholding (DAAP) Method
\begin{frame}{Gating Strategies: Detailed Thresholding (DAAP)}
The thresholding method is based on **Density-Adaptive Activation Pruning (DAAP)**:
\begin{itemize}
    \item **Formula**: Gating threshold $\tau_f = \tau_{\text{base}} / \max(\rho_f, 0.01)$, where $\rho_f$ is the non-zero weight density of filter $f$.
    \item **Methodology**: Sparse filters get higher thresholds (more selective gating); dense filters get lower thresholds.
    \item **Calibration**: Optimal base threshold ($\tau_{\text{base}}$) is found offline by sweeping and selecting the maximum value within a $0.5\%$ accuracy drop tolerance:
    \begin{itemize}
        \item **MNIST**: $\tau_{\text{base}} = 0.30$, yielding a threshold range of $0.30$--$0.34$.
        \item **CIFAR-10**: $\tau_{\text{base}} = 0.003$, yielding a threshold range of $0.0051$--$0.0069$.
    \end{itemize}
    \item **Hardware**: Implemented using zero-latency combinational comparators and distributed RAM threshold storage.
\end{itemize}
\end{frame}

% Slide 4: Lookahead Step Length Logic
\begin{frame}{Lookahead Step Gating Logic}
We implement lookahead logic to skip cycles when activations are masked:
\begin{itemize}
    \item **FC Layers**: Use a **4-step lookahead** ($L=4$). Addressing is purely linear ($input\_idx + 1$, $+2$, $+3$, etc.), keeping the critical path short and closing timing.
    \item **Conv Layers**: Restricted to a **2-step lookahead** ($L=2$). 
\end{itemize}
\begin{block}{Why 4-step lookahead was NOT used in Conv Layers}
\begin{enumerate}
    \item **Address Complexity**: Conv addressing requires coordinate calculations (kernel rows/columns, input channels, border wrapping).
    \item **Timing Violations**: Resolving a 4-step lookahead priority encoder with coordinate wrapping combinationally creates a path with $>50$ logic levels, violating timing and dropping $F_{\text{max}}$.
\end{enumerate}
\end{block}
\end{frame}

% Slide 5: Quantitative Performance Summary
\begin{frame}{Performance Summary: MNIST \& CIFAR-10}
\begin{table}
\centering
\resizebox{\textwidth}{!}{
\begin{tabular}{llcccc}
\toprule
\textbf{Dataset} & \textbf{Model Variant} & \textbf{Accuracy} & \textbf{Latency} & \textbf{LUTs} & \textbf{DSPs} \\
\midrule
MNIST & Baseline (Sequential) & 98.35\% & 0.832 ms & 14,925 (28\%) & 23 (10\%) \\
MNIST & TTQ + BN & 97.28\% & 0.905 ms & 14,588 (27\%) & 51 (23\%) \\
MNIST & TTQ + Threshold & 97.25\% & 0.887 ms & 23,085 (43\%) & 54 (24\%) \\
MNIST & TTQ + Hysteresis & 96.96\% & 0.815 ms & 33,565 (63\%) & 54 (24\%) \\
\midrule
CIFAR-10 & Baseline (KD Student) & 90.62\% & 70.50 ms & 23,241 (43\%) & 88 (40\%) \\
CIFAR-10 & TTQ + BN & 90.62\% & 71.94 ms & 24,689 (46\%) & 36 (16\%) \\
CIFAR-10 & TTQ + Threshold & 89.85\% & 50.45 ms & 31,828 (60\%) & 36 (16\%) \\
CIFAR-10 & TTQ + Hysteresis & 88.34\% & 36.88 ms & 35,798 (67\%) & 36 (16\%) \\
\bottomrule
\end{tabular}
}
\end{table}
\begin{itemize}
    \item For CIFAR-10, TTQ+BN retains the **full $90.62\%$ baseline accuracy**.
    \item Thresholding and Hysteresis yield speedups (down to 36.88 ms, a **$1.95\times$ speedup**) with minimal accuracy degradation ($89.85\%$ and $88.34\%$).
\end{itemize}
\end{frame}

% Slide 6: Industrial Use Cases & Advantages
\begin{frame}{Industrial Use Cases and Quantization Advantages}
\begin{block}{Why Quantization is Used}
\begin{itemize}
    \item **16$\times$ Weight Compression**: Compressing weights from FP32 to 2-bit ternary codes allows the network parameters to fit entirely on-chip, eliminating slow and power-hungry external DDR DRAM accesses.
    \item **No DSP Multipliers**: Ternary weights replace MAC multipliers with simple adders/subtractors, saving valuable silicon area.
    \item **Energy Efficiency**: Low dynamic power dissipation (~160 mW) prevents thermal throttling, keeping junction temp rise $<2^\circ$C.
\end{itemize}
\end{block}

\begin{block}{Autonomous Vehicles (AVs) Application}
\begin{itemize}
    \item **Latency-Critical Decision Loops**: AV systems require ultra-low latencies ($<5$ ms) for pedestrian and obstacle classification. Skip-gating FSMs dynamically speed up processing by skipping background pixels.
    \item **Passive Cooling**: Low-power FPGAs remove the need for bulky liquid cooling systems, reducing vehicle weight and battery drain.
\end{itemize}
\end{block}
\end{frame}

% Slide 7: Key Learnings & Takeaways
\begin{frame}{Key Learnings & Takeaways}
\begin{enumerate}
    \item \textbf{Quantization vs. Resource Trade-off}:
    \begin{itemize}
        \item Asymmetric Wp/Wn scale factors in TTQ successfully retain classification accuracy while reducing weight memory footprint.
    \end{itemize}
    \item \textbf{Data-Dependent Gating Success}:
    \begin{itemize}
        \item Spatial hysteresis is effective for sparse inputs (strokes, boundaries, text) but fails on dense, natural textures.
    \end{itemize}
    \item \textbf{Hardware Optimization}:
    \begin{itemize}
        \item Memory partitioning is essential to prevent BRAM port contention and avoid compiling weights into LUTs.
        \item Ping-pong buffering saves up to 7 BRAM36 tiles, proving the value of FSM schedule-aware routing.
    \end{itemize}
\end{enumerate}
\end{frame}

\end{document}
